\section{Actual adjoints and analytic conventions}\label{sec:conventions}

We collect the conventions that keep the argument on the given
holomorphic line bundles.  All complex spaces are Hausdorff and countable
at infinity.  A normal connected complex space is irreducible: its locally
irreducible components are open as well as closed.  All divisors called
effective are Weil divisors with nonnegative coefficients.

\subsection{Adjoints, residues, and singularities}

For a normal complex space \(X\), let \(j:X_{\mathrm{reg}}\hookrightarrow X\)
be the regular locus and set \(\omega_X=j_*\omega_{X_{\mathrm{reg}}}\).
For an integral Weil divisor \(H\), \(\OO_X(H)\) is the rank-one reflexive
divisorial sheaf.  If \(\Delta\) is rational and an integer \(r>0\)
clears its coefficients, the adjoint index condition is the invertibility
of
\begin{equation}
 \mathscr L_r=\bigl(\omega_X^{[r]}\otimes\OO_X(r\Delta)\bigr)^{**}.
 \label{eq:actual-adjoint}
\end{equation}
All later degrees are chosen divisible by such an index.  Tensor powers
of \(\mathscr L_r\), rather than a separately chosen global canonical
divisor, define \(\OO_X(m(K_X+\Delta))\).  Equality in \(\operatorname{Pic}(X)_\Q\)
means an isomorphism of holomorphic line bundles after a common positive
multiple.  The notation \(P\sim_\Q K_X+\Delta\) for a rational
\(\Q\)-Cartier divisor means precisely
\(\OO_X(mP)\simeq\OO_X(m(K_X+\Delta))\) in such a degree.

A local frame of \(\mathscr L_r\) restricts on the regular locus to a
meromorphic \(r\)-pluricanonical form with the allowed boundary poles.
Pulling that form back differentially along a proper bimeromorphic model
\(p:W\to X\) with \(W\) smooth defines a rational crepant subboundary \(G_W\) by
\begin{equation}
 K_W+G_W=p^*(K_X+\Delta).
 \label{eq:crepant}
\end{equation}
This equality records both the rational divisor of the meromorphic map
and the actual isomorphism of the pulled-back line.  It is independent
of the local frame: a change of frame is a holomorphic unit, and both
pullbacks change by its pullback.  At a prime divisor \(E\) on \(W\),
\[
 a(E;X,\Delta)=1-\operatorname{coeff}_E(G_W)
 =1+\operatorname{coeff}_E\bigl(K_W-p^*(K_X+\Delta)\bigr)
\]
is the log discrepancy.  We use the same definition on higher smooth
models.  The pair is lc if all these numbers are nonnegative and klt if
they are positive.  We use the standard Koll\'ar--Mori dlt notion, as
imported for analytic pairs in
\cite[Definition~2.7(4)]{DasHaconPaun2024}; in particular, a dlt pair is
klt away from its coefficient-one floor.  Theorem~\ref{thm:floor-generation}
uses in addition the globally simple projective resolution supplied by
Lemma~\ref{lem:special-resolution} in the fourfold application.
Theorem~\ref{thm:supported-lifting} uses the standard dlt notion without
this additional resolution hypothesis.

On an SNC pair, adjunction to an intersection of distinct coefficient-one
components is the iterated meromorphic residue.  In a degree divisible
by the adjoint index and by two, interchanging two residue operations
does not change the resulting pluriform.  We always use such even degrees
when comparing paths through boundary strata.  Equality of the pulled-back
invertible subsheaves of meromorphic pluriforms is stronger than an abstract
isomorphism of lines: the former fixes the residue transport of sections.
This stronger equality is what we mean by a \emph{crepant residue
comparison}.  It will be proved for every comparison used below.

Reflexive extension is used in the following precise form.  An isomorphism
between rank-one reflexive sheaves on a normal space, defined away from an
analytic subset of codimension at least two, extends uniquely if it is
given there by the same meromorphic identification.  In particular, a
meromorphic pluriadjoint identity transported through a bimeromorphic map
that extracts no prime divisor can be tested in codimension one and then
extended.  An isomorphism merely in numerical cohomology would not support
this operation.

We use ordinary global \(\Q\)-factoriality: every global Weil divisor has
a Cartier multiple, and a reflexive power of the canonical sheaf is
invertible.  This is the convention of
\cite[Definition~2.7]{DasHaconPaun2024}.  It does not assert factoriality
of all analytic germs.  The input of Theorem~\ref{thm:main} has no such
factoriality assumption.

\subsection{K\"ahler positivity and section spaces}

A K\"ahler form on a complex space is given by smooth strictly
plurisubharmonic local potentials in local embeddings into complex
Euclidean spaces.  Smooth metrics and their curvature are interpreted
in the same way.  For a rational line represented by \(L_r\), we use the
metric criterion for analytic nefness stated in the introduction.  It is
independent of the index and of the fixed K\"ahler form.  Restriction to a
closed analytic subspace preserves the inequalities and hence nefness.
For a holomorphic map from a compact K\"ahler space, pullback also preserves
nefness: after pulling back the metric inequality, bound the pulled-back
fixed form by a constant multiple of a fixed K\"ahler form on the source.
This argument will be used for resolutions and for restrictions to
possibly singular strata.  Curve intersection tests occur only in the
relative projective MMP calculations where their use is justified.

For an actual rational line \(D\), its Iitaka dimension is \(-\infty\)
if all positive Cartier powers have no sections; otherwise it is the
maximum dimension of the meromorphic images of their complete systems.
On an irreducible compact space, two linearly independent sections of one
line have a nonconstant meromorphic quotient.  Consequently, if
\(\kappa(D)=0\), every Cartier degree has at most one independent section.
Conversely, unbounded dimensions of the section spaces of positive powers
force \(\kappa(D)\geq1\): one of those spaces has two independent sections.

We will repeatedly use two elementary consequences of normality.  If
\(p:Y\to X\) is a proper bimeromorphic morphism between normal spaces, then
\(p_*\OO_Y=\OO_X\).  Hence, for a line \(L\) on \(X\),
\[
 H^0(Y,p^*L)=H^0(X,L).
\]
If \(p^*L\) is generated, then \(L\) is generated: at any \(x\in X\)
choose \(y\in p^{-1}(x)\) and a pulled-back section nonzero at \(y\).
The corresponding section of \(L\) is nonzero at \(x\).  If a positive
multiple of \(D\) has a nowhere-vanishing section, that section is an
actual isomorphism \(\OO_X\simeq\OO_X(mD)\).

The second consequence concerns exceptional divisors.  If \(E\geq0\)
is \(p\)-exceptional and integral, then
\begin{equation}
 p_*\OO_Y(E)=\OO_X. \label{eq:exceptional-hartogs}
\end{equation}
Indeed, a section is a meromorphic function on \(X\) that is holomorphic
away from the codimension-at-least-two image of \(E\); normal Hartogs
extension makes it holomorphic everywhere.  The same statement applies
to rational exceptional divisors after clearing their denominators.
Combined with projection formula, it identifies the section spaces in
the exceptional comparisons below.

\subsection{A connectedness and rationality lemma}

The following sheaf calculation is used both for the conormal layers of
a prime floor and for descent of compatible boundary sections.  Its klt
case also supplies rationality and the Cohen--Macaulay property for the
strata used later.

\begin{lemma}\label{lem:floor-connectedness}
Let \((T,B_T)\) be a normal irreducible effective lc analytic pair with actual
\(\Q\)-Cartier adjoint, and let \(p:U\to T\) be a projective resolution
with smooth source whose crepant subboundary \(G\) has SNC support.  Put
\[
 R=G^{=1},\qquad H=\ceil{-G^{<1}}.
\]
Then \(H\) is effective and exceptional, has no component in common with
\(R\), and
\begin{equation}\label{eq:floor-pushforward}
 p_*\OO_R=\OO_{p(R),\red}.
\end{equation}
In particular \(R\to p(R)\) has connected fibers.  If the pair is klt,
then \(T\) has rational singularities and is Cohen--Macaulay.
\end{lemma}
\begin{proof}
The strict boundary coefficients belong to \([0,1]\), so positive
coefficients of \(H\) occur only on exceptional primes.  Coefficient by
coefficient,
\[
 H-R=-\floor G,
 \qquad (H-R)-(K_U+\{G\})=-p^*(K_T+B_T).
\]
The fractional boundary \(\{G\}=G-\floor G\) is effective SNC with
coefficients below one.  The displayed adjoint difference is relatively
nef and big for the bimeromorphic \(p\): it is relatively trivial, and
the generic relative dimension is zero.  Analytic relative
Kawamata--Viehweg vanishing
\cite[Theorem~2.41]{DasHaconPaun2024} gives
\(R^ip_*\OO_U(H-R)=0\) for \(i>0\).  Exceptional Hartogs extension gives
\(p_*\OO_U(H)=\OO_T\).  The sequence for the Cartier divisor \(R\)
therefore yields a surjection
\[
 \OO_T\longrightarrow p_*\OO_R(H).
\]
It factors through \(p_*\OO_R\).  Multiplication by the canonical section
of \(H\) embeds \(\OO_R\) into \(\OO_R(H)\), because \(R\) is reduced
SNC and shares no component with \(H\).  The surjection consequently
makes that embedding an isomorphism after pushforward and makes
\(\OO_T\to p_*\OO_R\) surjective.  Its kernel is exactly the reduced
ideal of the proper image \(p(R)\): a holomorphic function vanishes after
restriction to \(R\) precisely when it vanishes on that image.  This
proves \eqref{eq:floor-pushforward}; Stein factorization gives connected
fibers.

For the last assertion, take a klt resolution, so \(R=0\).  The same
vanishing and Hartogs calculation identify
\(Rp_*\OO_U(H)\simeq\OO_T\).  The factorization
\[
 \OO_T\longrightarrow Rp_*\OO_U\longrightarrow Rp_*\OO_U(H)\simeq\OO_T
\]
is the identity, as can be checked on degree zero.  It splits the first
map.  Apply proper coherent duality.  Canonical higher direct images of
the resolution vanish: the canonical torsion-freeness theorem
\cite[Theorem~2.9 and Proposition~2.11]{Fujino2023AnalyticVanishing}
applies locally on the base, and these higher images vanish generically
for a bimeromorphic map.  The required local K\"ahlerness follows by
restricting to a small Stein base neighborhood and combining a relative
ample metric with a pulled-back strictly plurisubharmonic potential.
Thus the dual split surjection is the trace
\[
 p_*\omega_U[\dim U]\longrightarrow\omega_T^\bullet,
\]
where \(\omega_T^\bullet\) is the dualizing complex.  It follows that this
complex is concentrated in degree \(-\dim U\).  On that cohomology sheaf
the trace is also injective: its source is torsion-free and the map is
generically an isomorphism.  Hence it is an isomorphism.  Biduality gives
\(Rp_*\OO_U=\OO_T\).  This is rational singularity, and the concentration
of the dualizing complex is the Cohen--Macaulay property.
\end{proof}
